\section{Notation}

We here use the tensor network notation \cite{goessmann_tensor_2026}.
The modes of each tensors are assigned by variables, and denoted in brackets as
\begin{align*}
    \hypercorewith \, .
\end{align*}
Quantum gates are along this line unitary tensors with incoming and outgoing variables.
The Pauli-X for example is denoted by
\begin{align*}
    \paulixat{\cavariable{\insymbol},\cavariable{\outsymbol}} \coloneqq
    \coloredmatrixof{
        0 & 1 \\
        1 & 0}{\cavariable{\insymbol},\cavariable{\outsymbol}} \, .
\end{align*}
Contractions are denoted by angle brackets with the variables left open specified in further brackets.
Circuits are contractions of the unitaries, leaving only the last output variables open.
Computational basis measurements of circuits are coordinatewise squared absoluted of the contracted circuit.